\subsection{DELB-Based Evaluation Algorithm}
\label{subsec:delb_evaluation_algorithm}

\begin{Algorithm}[DELB-based evaluation of an added urban data source]
\label{alg:delb_evaluation}
\textbf{Input:} integer target \(C\); baseline state \(\mathcal S^0\);
added-data state \(B\); training, validation, and test indices; entropy
estimator; randomization designs; number of repetitions \(R\).\\
\textbf{Output:} \(\widehat L^0\), \(\widehat L^B\),
\(\widehat{\Delta L}\), \(\widehat I(C;B\mid\mathcal S^0)\), support
diagnostics, randomization \(p\)-values, and matched held-out prediction
errors.
\begin{enumerate}
    \item Estimate all discretization thresholds from the training period
    and freeze them for validation and testing.
    \item Construct the baseline state \(\mathcal S^0\) and augmented state
    \(\mathcal S^B=(\mathcal S^0,B)\) on identical eligible observations.
    \item Estimate \(\widehat H(C\mid\mathcal S^0)\) and
    \(\widehat H(C\mid\mathcal S^B)\) with the same entropy estimator.
    \item Numerically evaluate
    \(\widehat L^j=\mathcal H_{\mathbb Z}^{-1}
    [\widehat H(C\mid\mathcal S^j)]\), \(j\in\{0,B\}\).
    \item Compute raw CMI, the absolute and relative estimated DELB
    reductions, and the support-density, singleton-share, and effective
    degrees-of-freedom diagnostics.
    \item For each prespecified randomization design and repetition, transform
    \(B\), reconstruct \(\mathcal S^B\), and recompute both state support and
    the complete statistic; then calculate plus-one upper-tail \(p\)-values.
    \item Fit each prediction model twice with matched samples, training
    windows, tuning rules, and all non-bicycle inputs: once with
    \(\mathcal S^0\) and once with \(\mathcal S^B\).
    \item Compare the estimated DELB change, randomization-calibrated evidence,
    and held-out MSE change as three separate outputs.
\end{enumerate}
\end{Algorithm}

\paragraph{Numerical inversion.}
For a positive entropy value, the implementation brackets the
discrete-Gaussian rate \(\lambda\) from an initial value of \(1\) by repeated
halving or doubling and then applies Brent's method. The root tolerance is
\(10^{-12}\), with at most 200 iterations. For \(\lambda\geq1\), the
partition function and second moment are evaluated by direct symmetric
integer-lattice summation. For \(\lambda<1\), the Poisson-dual theta series
avoids an unnecessarily wide real-lattice truncation. In both cases the
support radius is expanded until the analytic bound on the omitted
two-sided unnormalized tail is at most \(10^{-15}\). Entropy values at or
below zero map to a DELB of zero. Repeated resampling and randomization paths
use the configured monotone interpolation grid over 0--8 bits, with direct
inverse evaluations retained for validation.

\paragraph{Computational cost and parallelization.}
Let \(N\) be the number of matched observations, \(K_0\) and \(K_B\) the
numbers of occupied atoms in the two state tables, \(M\) the maximum retained
half-support of a lattice sum, \(T\leq200\) the number of root evaluations,
and \(J\) the number of entropy values converted to bounds. State counting
requires \(O(N)\) time and \(O(K_0+K_B)\) memory; direct bound conversion
requires \(O(JTM)\) time and \(O(M)\) additional memory. With \(R\)
randomizations, the serial calibration cost is
\(O(RN+RJTM)\), excluding model fitting, whose cost depends on the selected
predictor. Randomizations are independent conditional on their stored seeds.
The implementation evaluates city--spatial-scale--temporal-scale cells
independently and records cell-level results and seeds for reproducibility.
